\setcounter{neje}{68}
\begin{nejer}[Imagen directa de una medida]
Sea $(X, \mathcal{A}, \mu)$ un espacio de medida, sea $Y$ un conjunto y sea $\phi: X \to Y$ una aplicación. La imagen directa (\textit{pushforward}) de $\mathcal{A}$ es la $\sigma$-álgebra
$$
\phi_*\mathcal{A} := \set*{B \subseteq Y \mid \phi^{-1}(B) \en \mathcal{A}} \subseteq 2^Y.
$$
La imagen directa de $\mu$ es la función
$$
\phi_*\mu\colon \phi_*\mathcal{A} \to [0, \infty]
$$
definida por $(\phi_*\mu)(B) := \mu(\phi^{-1}(B))$ para $B \en \phi_*\mathcal{A}$.

\begin{enumerate}[label = (\roman*)]
\item Pruebe que $(Y, \phi_*\mathcal{A}, \phi_*\mu)$ es un espacio de medida.

\item Sea $(X, \mathcal{A}^*, \mu^*)$ la completación de $(X, \mathcal{A}, \mu)$ y sea $(Y, (\phi_*\mathcal{A})^*, (\phi_*\mu)^*)$ la completación de $(Y, \phi_*\mathcal{A}, \phi_*\mu)$. Pruebe que
$$
(\phi_*\mu)^*(E) = \mu^*(\phi^{-1}(E)), \quad \forall E \en (\phi_*\mathcal{A})^* \subseteq \phi_*\mathcal{A}^*.
$$
Deduzca que $(Y, \phi_*\mathcal{A}, \phi_*\mu)$ es completo siempre que $(X, \mathcal{A}, \mu)$ sea completo. Encuentre un ejemplo donde $(\phi_*\mathcal{A})^* \subsetneq \phi_*\mathcal{A}^*$.

\item Fije una función $f\colon Y \to [0, \infty]$. Pruebe que $f$ es $\phi_*\mathcal{A}$-medible si y solo si $f \circ \phi$ es $\mathcal{A}$-medible. Si $f$ es $\phi_*\mathcal{A}$-medible, pruebe que
$$
\int_Y f \dd{(\phi_*\mu)} = \int_X (f \circ \phi) \dd{\mu}.
$$

\item Determine la imagen directa de $(X, \mathcal{A}, \mu)$ bajo una aplicación constante.
\end{enumerate}
\end{nejer}
\begin{sol}\vphantom{r}
\begin{enumerate}[label = (\roman*)]
%-----------------------------------------------------
\item %(i)
\emph{Veamos que $\phi_*\mathcal{A}$ es $\sigma$-álgebra}:\\
Por un lado, $\varnothing \en \mathcal{A}$ garantiza que $\varnothing=\phi^{-1}(\varnothing)\en\phi_*\mathcal{A}$. Por otro lado, si $\set{A_n}_{n\en\mathbb{N}}$ es una familia de elementos de $\phi_*\mathcal{A}$, entonces $\set*{\phi^{-1}(A_n)}_{n\en\mathbb{N}}$ es una familia de elementos de $\mathcal{A}$, por lo que $\phi^{-1}(\bigcup^\infty_{n=1}A_n) = \bigcup^\infty_{n=1}\phi^{-1}(A_n) \en \mathcal{A}$, lo cual implica que $\bigcup^\infty_{n=1}A_n\en\phi_*\mathcal{A}$. Finalmente, si $A\en\phi_*\mathcal{A}$ entonces $\phi^{-1}(A) \en \mathcal{A}$, de donde $\phi^{-1}(Y\setminus A) = X\setminus \phi^{-1}(A) \en \mathcal{A}$, por lo que $Y\setminus A \en \phi_*\mathcal{A}$.\vspace{1em}\\
\emph{Veamos que $\phi_*\mu$ es medida}:\\
Primero notemos que $\phi_*\mu(\varnothing) = \mu(\phi^{-1}(\varnothing)) = \mu(\varnothing) = 0$. Además, si $\set{A_n}_{n\en\mathbb{N}}$ es una familia de elementos de $\phi_*\mathcal{A}$ disjuntos a pares, se sigue que $\set{\phi^{-1}(A_n)}_{n\en\mathbb{N}}$ es una familia de elementos de $\mathcal{A}$, los cuales también serán disjuntos a pares gracias a que la imagen inversa preserva intersecciones. Luego,
\[
\phi_*\mu\biggl(\, \bigcup^\infty_{n=1}A_n \,\biggr)
= \mu\biggl(\phi^{-1}\biggl(\, \bigcup^\infty_{n=1}A_n \,\biggr)\biggr)
= \mu\biggl(\, \bigcup^\infty_{n=1} \phi^{-1}(A_n) \,\biggr)
= \sum^\infty_{n=1}\mu(\phi^{-1}(A_n))
= \sum^\infty_{n=1} \phi_*\mu(A_n)
\]
%-----------------------------------------------------
\item %(ii)
Sea $E \en (\phi_*\mathcal{A})^*$. Por definición existen $A,B\en\phi_*\mathcal{A}$ con $A \subseteq E \subseteq B$ y $\phi_*\mu(B \setminus A) = 0$; más aún, también se define $(\phi_*\mu)^*(E) = \phi_*\mu(A)$. De esto se sigue que $\phi^{-1}(A) \subseteq \phi^{-1}(E) \subseteq \phi^{-1}(B)$ y $\mu(\phi^{-1}(B) \setminus \phi^{-1}(A)) = \mu(\phi^{-1}(B\setminus A)) = \phi_*\mu(B\setminus A) = 0$, como $\phi^{-1}(A),\phi^{-1}(B) \en \mathcal{A}$, esto demuestra que $\phi^{-1}(E) \en \mathcal{A}^*$, de donde $E \en \phi_*\mathcal{A}^*$; más aún, se define $\mu^*(\phi^{-1}(E)) = \mu(\phi^{-1}(A))$. Pero ya vimos que $(\phi_*\mu)^*(E) = \phi_*\mu(A)$, siendo esto último igual a $\mu(\phi^{-1}(A))$. Por lo tanto, $(\phi_*\mu)^*(E) = \mu^*(\phi^{-1}(E))$.\par
De esto se sigue que, si $(X,\mathcal{A},\mu)$ es completo, $(\phi_*\mathcal{A})^* \subseteq \phi_*\mathcal{A}^* = \phi_*\mathcal{A}$ (ya que $\mathcal{A} = \mathcal{A}^*$ gracias al \myref{nejer:1.67}). Esto implica que $\phi_*\mathcal{A}$ es completo: si $A \subseteq F$ con $F \en\phi_*\mathcal{A}$ y $\phi_*\mu(F) = 0$, entonces $\varnothing \subseteq A \subseteq F$ con $\phi_*\mu(F \setminus\varnothing) = 0$, así que $A \en (\phi_*\mathcal{A})^* \subseteq \phi_*\mathcal{A}$, de donde $A \en \phi_*\mathcal{A}$.
\vspace{1em}\\
\emph{Encontremos un ejemplo de la contención propia}:\\
Sea $(X,\mathcal{A},\mu)$ cualquier espacio de medida incompleto con $\mu(X)>0$,%
\footnote{
    Un ejemplo clásico puede ser $(\mathbb{R},\mathcal{B}_{\mathbb{R}},\lambda)$. Un ejemplo minimal es $X = \set{0,1,2}$, $\mathcal{A} = \set[\big]{\varnothing, \set0, \set{1,2}, X}$ y $\mu = \delta_0$ la medida de Dirac concentrada en $0$. Este espacio es incompleto ya que $\set1 \subseteq\set{1,2}$ con $\delta_0(\set{1,2}) = 0$, pero $\set1 \noten \mathcal{A}$.
}
así que existe $N\noten \mathcal{A}$ tal que $N \subseteq B$ con $B\en\mathcal{A}$ y $\mu(B)=0$ (de esto se sigue que $N\en\mathcal{A}^*$ y $\mu^*(N) = 0$). También tomemos $Y = \set{0,1}$ y $\phi \colon X \to Y$ como el mapeo característico de $N$, i.e., $\phi = \chi_N$.\par
Ahora bien, como $\phi^{-1}(\set1) = N$, tenemos que $\phi^{-1}(\set1) \en \mathcal{A}^*$ y $\phi^{-1}(\set1)\noten \mathcal{A}$, es decir, $\set1 \en \phi_*\mathcal{A}^*$ y $\set1\noten\phi_*\mathcal{A}$. Solo resta probar que $\set1\noten(\phi_*\mathcal{A})^*$:\par
En efecto, si suponemos por contradicción que $\set1\en(\phi_*\mathcal{A})^*$, entonces existen $E,F\en\phi_*\mathcal{A}$ tales que $E \subseteq \set1 \subseteq F$ y $\phi_*\mu(F\setminus E) = 0$; en particular, se cumple $(\phi_*\mu)^*(\set1) = \phi_*\mu(F)$. Sin embargo, recordemos que $Y = \set{0,1}$, así que $\set1 \subseteq F \subseteq \set{0,1}$; en otras palabras, $F = \set1$ o $F=\set{0,1}$, pero $F\en\phi_*\mathcal{A}$ y ya probamos que $\set1\noten\phi_*\mathcal{A}$, por lo que forzosamente $F = \set{0,1}$. Esto garantiza que 
\[
(\phi_*\mu)^*(\set1) = \phi_*\mu(F) =\mu(\phi^{-1}(\set{0,1})) = \mu(X) > 0
\]
pero por lo anteriormente probado en (ii): $(\phi_*\mu)^*(\set1) = \mu^*(\phi^{-1}(\set1)) = \mu^*(N) = 0$. Esto es una contradicción. El error estuvo en suponer que $\set1\en(\phi_*\mathcal{A})^*$, así que $\set1\noten(\phi_*\mathcal{A})^*$, y queda probado que $\set1\en\phi_*\mathcal{A}^*\setminus(\phi_*\mathcal{A})^*$, por lo que la contención es propia.
%-----------------------------------------------------
\item %(iii)
Por definición, $f$ es medible si y solo si $f^{-1}(A) \en \phi_*\mathcal{A}$ para todo boreliano $A$, lo cual equivale por definición que $\phi^{-1}(f^{-1}(A))\en\mathcal{A}$ para todo boreliano $A$. Como $(f\circ\phi)^{-1}(A) = \phi^{-1}(f^{-1}(A))$, esto prueba directamente que $f$ es $\phi_*\mathcal{A}$-medible si y solo si $f\circ\phi$ es $\mathcal{A}$-medible.\par
Ahora bien, si $f$ es $\phi_*\mathcal{A}$-medible, existe $\set{s_n}_{n}$ sucesión monótona de funciones escalonadas, medibles y no negativas que converge puntualmente a $f$. Luego, por convergencia monótona
\begin{equation}\label{eq:aprox1}
\lim_{n\to\infty}\int_Y s_n \dd{(\phi_*\mu)} = \int_Y f \dd{(\phi_*\mu)}
\end{equation}
Como $\im(s_n \circ \phi) \subseteq \im(s_n)$, se sigue que $s_n\circ\phi$ también es escalonada no negativa, y será medible gracias a que $s_n$ lo es. Más aún, como $\set{s_n}_n$ converge monótona y puntualmente a $f$ entonces se sigue, para cada $x\en X$, que $s_{n+1}(\phi(x)) \ge s_n(\phi(x))$ y $\lim_{n\to\infty} s_n(\phi(x)) = f(\phi(x))$, así que $\set{s_n\circ\phi}_n$ converge monótona y puntualmente a $f \circ \phi$. Por lo tanto, la convergencia monótona garantiza que 
\begin{equation}\label{eq:aprox2}
\lim_{n\to\infty}\int_X s_n\circ\phi \dd{\mu} = \int_X f\circ\phi \dd{\mu}
\end{equation}
Sin embargo, si ponemos $s_n = \sum^{m}_{k=1}a_k\chi_{A_k}$, entonces $s_n\circ\phi = \sum^m_{k=1}a_k\chi_{\phi^{-1}(A_k)}$, de donde
\[
\int_Y s_n \dd{(\phi_*\mu)}
= \sum^m_{k=1} a_k \cdot\phi_*\mu(A_k)
= \sum^m_{k=1} a_k \cdot\mu(\phi^{-1}(A_k))
= \int_X s_n \circ \phi \dd{\mu}
\]
Como esto se cumple para cualquier $n\en\mathbb{N}$, (\ref{eq:aprox1}) y (\ref{eq:aprox2}) garantizan que
\[
\int_Y f \dd{(\phi_*\mu)} = \int_X f\circ\phi \dd{\mu}
\]
%-----------------------------------------------------
\item %(iv)
Si $\phi\colon X\to Y$ es de valor constante $y_0\en Y$, entonces para cualquier $A\subseteq Y$ se tiene:
\[
\phi^{-1}(A)
= 
\begin{cases}
    X,           &\text{si }y_0\en A.\\
    \varnothing, &\text{si }y_0\noten A.
\end{cases}
\]
En cualquier caso, se sigue que $\phi^{-1}(A)\en\mathcal{A}$, por lo cual $\phi_*\mathcal{A} = 2^Y$. Más aún, lo anterior también implica que para todo $A \subseteq Y$ se tiene:
\[
\phi_*\mu(A) = \mu(\phi^{-1}(A))
=
\begin{cases}
    \mu(X), &\text{si } y_0 \en A.\\
    0,      &\text{si } y_0 \noten A.
\end{cases}
\]
Esto prueba que $\phi_*\mu = \mu(X)\cdot\delta_{y_0}$ es la medida de Dirac concentrada en $y_0$ multiplicada por $\mu(X)$. Hemos probado que $(Y,\phi_*\mathcal{A},\phi_*\mu) = (Y,2^Y,\mu(X)\delta_{y_0})$.
\end{enumerate}
\end{sol}